\subsection{原始三乗根の量子三角形の全中央切断}
\label{subsec:cubic-CY-three-alpha-primitive-triangle}

$q$ を原始三乗根とする。以下では
命題\ref{prop:cubic-CY-three-alpha-quantum-triangle-incidence}の
$A_i,B_i,X,Y,Z,\Omega,t$ を用いる。
固定出典の三角形の定理により、六係数の全値でこの包絡はKoszul AS4かつCYである
\cite{AS4CYpinned}。
活動次数濾過の随伴環は $t$ を加えた量子多項式環なので、
包絡はnoetherian整域である。

\begin{proposition}\label{prop:cubic-CY-three-alpha-primitive-triangle-centre}
$T=t^3$ とし、巡回三組 $(i,j,k)$ を $(x,y,z),(y,z,x),(z,x,y)$ とする。
$X_i$ をそれぞれ $X,Y,Z$ と書き、$d=(2q+1)/9$、
\[
\begin{aligned}
 L_i={}&-\frac{2(q-1)}9A_i^2A_jA_k B_i
         -\frac q3 A_i^2B_jB_k-dB_i^3,\\
 L_\Omega={}&-\frac{q+1}3\sum_i A_jA_kB_i^2-\frac{q+2}3B_xB_yB_z,\\
 C_0={}&-\frac{q+2}9\sum_i A_iA_j^2A_k^2B_i^2
         +\frac{q-1}9\sum_i A_iB_j^2B_k^2
\end{aligned}
\]
と置く。全中心は $t,X,Y,Z,\Omega$ の生成する可換環で、
唯一の関係は次の式である。
\begin{equation}\label{eq:cubic-CY-three-alpha-primitive-triangle-F}
\begin{aligned}
 \Omega^3&=XYZ+
 T\left(\frac{q-1}3 A_xA_yA_z\Omega^2+
          \Omega\sum_i A_iB_iX_i+d\sum_i A_i^3X_i^2\right)\\
 &\quad+T^2\left(L_\Omega\Omega+\sum_i L_iX_i\right)+T^3 C_0 .
\end{aligned}
\end{equation}
包絡は $k[t,X,Y,Z]$ 上で自由階数27である。
$t=1$ の包絡の中心も同じ式の $T=1$ で与えられ、
正規整域である。中心の分数体上で包絡は次数3の中心斜体となる。
\end{proposition}
\begin{proof}
活動次数の先頭で中央になる語の指数は全て同じ剰余を持つ。
従って最高項を $X,Y,Z,\Omega$ の積で引く帰納法により中心を生成できる。
この議論は $t=1$ でも同じである。
余り指数が各々 $0,1,2$ に入る順序語を選ぶと、
$k[t,X,Y,Z]$ 上で張り、最高語の指数の剰余により独立である。
中心では $1,\Omega,\Omega^2$ だけが同じ剰余の三類に対応し、
これらは $k[t,X,Y,Z]$ 上で独立である。
$\Omega^3$ を順序語へ還元して得る関係を $t$ で斉次化すると
\eqref{eq:cubic-CY-three-alpha-primitive-triangle-F} となる。
\path{verify_cubic_CY_three_alpha_primitive_triangle_centre.py}
は全六係数でこの恒等式を再計算する。

$t=1$ の中心は整域であり、上式を左辺へ集めた多項式を $F$ と書く。
その次数三の成分は $\Omega^3-XYZ$。
四偏導関数の最高成分は $3\Omega^2,-YZ,-XZ,-XY$ である。
この四元の零集合は三つの軸の和で次元1なので、特異集合の次元も高々1。
三次元の超曲面はCohen--Macaulayであり、Serreの判定から正規である。
包絡の中心上の階数は $27/3=9$。
中心の分数体への局所化は有限次元整域であり、従って中心斜体である。
\end{proof}

\begin{theorem}\label{thm:cubic-CY-three-alpha-primitive-triangle-test}
$G$ を $X,Y,Z,\Omega,t^3$ の非零線型結合とし、
少なくとも一つの活動係数を非零とする。
$R=\mathcal B/(G)$、$V=(\mathcal B[t^{-1}])_0$、$U=V/(G|_{t=1})$ と置く。
$C=Z(V)/(G|_{t=1})$ を上の三次式と切断の一次式から構成する。
次の条件がH3の必要十分条件である。
\begin{enumerate}
\item $C$ は整域で、有限 $C$ 加群 $U$ はtorsion-freeである。
\item $K=\operatorname{Frac}C$ 上の九次元代数 $U_K$ の根基は零であり、
中心の次元は1である。滑らかな射影モデル上で、
この次数3の中心単純代数のBrauer類に非零の剰余を持つ素因子がある。
\item $P=k[T,X,Y,Z]$ の四つの標準アフィンチャートで
\[
 A_f=\bigl(R^{(3)}[f^{-1}]\bigr)_0,
 \qquad f\in\{T,X,Y,Z\},
\]
と置く。それぞれの任意の二つの単純加群 $S_1,S_2$ に対し
$\operatorname{Ext}^3_{A_f}(S_1,S_2)=0$。
空のチャートではこの条件は自動的に成立する。
\end{enumerate}
第三条件の単純加群の次元は高々27である。
三条件はいずれも有限の厳密な代数手順で調べられる。
この判定は切断の $t=0$ の曲線の既約性や滑らかさを仮定しない。
\end{theorem}
\begin{proof}
\medskip\noindent\textbf{整域性}\par
$t=0$ で $G$ の像は零ではないので、$t$ は $R$ 上で正則である。
実際、$th=Gr$ なら $\mathcal B/(t)$ の整域性から $r=tr'$ となり、
$t$ を消して $h=Gr'$ を得る。
従って $R$ の整域性は $U$ の整域性と同値。
定数だけの切断は $t^3$ の倍数で整域にはならない。

$\overline G=G|_{t=1}$ とする。
$Z(V)\cap\overline G V=\overline G Z(V)$ である。
右辺への包含は、もし $\overline G a$ が中央であれば整域 $V$ で
$\overline G$ を消して $a$ の中央性を得ることによる。
従って $C\to U$ は単射であり、第一条件は整域性に必要である。
$C$ が整域なら $(\overline G)$ は $Z(V)$ の高さ1の素イデアル。
命題\ref{prop:cubic-CY-three-alpha-primitive-triangle-centre}の正規性から、
その局所環は離散付値環となる。$V$ はこの局所環上で有限torsion-free、
従って自由階数9である。剰余を取ると $\dim_K U_K=9$ となる。
第一条件の下では $U\to U_K$ は単射なので、
整域性は $U_K$ が斜体であることに同値である。

$U_K$ の元は全て $K$ 上のモニック三次式を満たす。
実際、離散付値環上の自由orderの元はその分数体でのreduced characteristic polynomialを満たす。
その係数は分数体に属し、通常の左乗法の特性多項式の根の対称式として整なので、
離散付値環に属する。これを剰余体に特殊化する。
九次元の斜体が中心拡大を持つなら、次元公式より九次の可換体拡大である。
標数零の原始元定理によりモニック三次条件に反する。
従って必要な斜体は中心が正確に $K$ であり、その次数は3。
逆に根基零と中心次元1は次数3の中心単純性を与える。
この素次数では斜体性はBrauer類の非零性と同値である。

$K$ は可換三次曲面の関数体である。
実際、$\overline G$ で一つの中心座標を消去し、得られる式の射影閉包を取る。
最高次数の $\Omega^3-XYZ$ は一次因子を持たず、その次数は消去後も3である。
整な三次曲面で二重点があれば、その点からの射影で有理となる。
三重点があれば平面三次曲線上の錐なので、その曲線の正規化上にruledとなる。
滑らかな三次曲面は代数閉体上で有理である。
従って $K$ の滑らかな射影モデルは有理曲面または滑らかな曲線上のruled曲面と
双有理であり、Brauer群は零となる。最後は曲面のBrauer群の双有理不変性、
射影束の公式およびTsenの定理による\cite{StacksTsenBrauer}。
滑らかな曲面上のpurityから、$\operatorname{Br}(K)$ の類が零である
必要十分条件は全素因子で剰余が零であること。
よって第二条件を得る。

\medskip\noindent\textbf{射影側の有限大域次元}\par
$\mathcal B^{(3)}$ は次数1の四中心元を持つ多項式環 $P$ 上で
自由階数27である。これは上の順序基底に $t$ の指数を $0,1,2$ の
剰余で一意に足して得られる。
$R^{(3)}$ はこれの有限商なので、$\qgr R$ は
$\mathbb P^3=\operatorname{Proj}P$ 上の対応する有限な代数層の
coherent加群圏と同値である。
標準四開集合での代数は $A_f$ であり、
各代数は三変数多項式環の上で高々27元で生成される。
各単純はこの中央多項式環の閉点ファイバーに属すので、次元は高々27。
系\ref{cor:cubic-CY-three-alpha-triangle-finite-effective}と同じ
半完全な完備化とTorの議論から、単純間の第三Extの消滅は
各 $A_f$ の大域次元が高々2であることと同値である。

第三条件が成り立てば局所Extと局所大域スペクトル系列から
$\qgr R$ の大域次元は有限である。
逆にH3であればSerre双対性から $\qgr R$ の第三Extは消える。
標準アフィンにセンターを局所化したExtは
中央式の冪に沿うtwist付きExtの直極限で計算できる。
coherent加群は標準アフィンから延長できるので、
各 $A_f$ で第三Extが消える。

整域切断では、$t=0$ の商は正次数でHilbert値 $3n$ を持つ。
この商は有限な射影代数層の加群であり、その構造対象は有限長ではない。
従って $t$ の降鎖とHilbert多項式から
Gabriel--Krull rank-one torsion-free性が従う。
残るH3公理は正則中央三次商の認識定理を適用して得る。

\medskip\noindent\textbf{有限判定}\par
$C$ の整域性は三次式の因数分解で判定する。
$U$ のtorsionは、一般点で選んだ基底の分母を払って得る
有限 $C$ 加群写像の核として、加群のGr\"obner基底で計算できる。
$U_K$ の根基と中心は、正標数を用いず、正則表現のtrace対形式と交換子の
線型方程式で判定できる。
次数3の中心単純代数はWedderburnの巡回性定理によりsymbol表示を持つ。
その表示の探索は、有理関数係数の次数上界を順に増やし、
三乗が非零スカラー、$q$ 交換関係、九基底語の独立性を
有限多項式条件として解くことで停止する。
symbol表示 $(f_0,g_0)_q$ を得れば、素因子の剰余は
$f_0^{v(g_0)}/g_0^{v(f_0)}$ の単元剰余の三乗類である。
調べる素因子は滑らかな射影モデル上の $f_0,g_0$ の零と極に限られる。
第三条件は各 $A_f$ の対角双加群の有限自由解消の最初の五項と、
次元27以下の単純表示の行列を用いる。
有限な多項式系の空性へ帰着する手順は
系\ref{cor:cubic-CY-three-alpha-triangle-finite-effective}と同じである。
\end{proof}

\begin{proposition}\label{prop:cubic-CY-three-alpha-primitive-triangle-arrows}
上の許容切断間の全点付き幾何学的矢は、共通スカラーを除き、
$t\mapsto d_0t'$、$x_i\mapsto\lambda_i x'_{\pi(i)}$ により与えられる。
ここで $d_0,\lambda_i\ne0$、$\pi\in S_3$ であり、
偶置換では $q'=q$, $\eta=1$、奇置換では $q'=q^{-1}$, $\eta=-q$ とする。
条件は
\[
 \eta\lambda_j\lambda_k A'_{\pi(i)}=d_0\lambda_i A_i,\qquad
 \eta\lambda_j\lambda_k B'_{\pi(i)}=d_0^2 B_i
\]
と、代入により切断直線 $kG$ が $kG'$ に移ることである。
特に慣性はこの明示作用での安定化群である。
\end{proposition}
\begin{proof}
中央一次の直線は $kt$ である。$t=0$ の三量子関係に各純平方が
現れないので、変換行列の各列で異なる二行の積が零になる。
従って変換行列は単項。$t$ に比例した平行移動は、
各関係で $t x'_j,t x'_k$ の係数を比較すると $q\ne1$ により消える。
元関係に代入すると表示した式となる。偶・奇の違いは
$q$ 交換子の語順の反転だけから生じる。
これらの式と切断直線の条件は逆に次数付き同型を与えるので、
反Serre環による内在的再構成から全矢となる。
\end{proof}
